% GENERATED FILE. Edit canonical lessons, then run npm run generate:books.
% canonical-source-hash: fnv1a64:fe55fd56d70e5d13
% canonical-lessons: GE-C169-bekannt, GE-C169-fremd, GE-C169-nervoes, GE-C169-stolz, GE-C169-wuetend

\chapter{Well known, Foreign, Nervous, Proud, Angry}
\label{ch:ge-well-known-foreign-nervous-proud-angry}
\hlchaptermodality{169}

\begin{chapteropening}
\textbf{By the end of this chapter:} \emph{I can say well known, foreign, nervous, proud and angry.}

Complete the last lesson of chapter 169: I can say well known, foreign, nervous, proud and angry.
\end{chapteropening}

\section[bekannt]{bekannt — well known}

\label{lesson:GE-C169-bekannt}



\begin{quote}
Before the new one: say the German for full, then the German for old-fashioned.
\end{quote}

\subsection*{You'll want to know: bekannt}
\textbf{bekannt} — \textquotedblleft{}well known\textquotedblright{}.

\begin{cousinweb}
\textbf{bekannt} is a past participle built on \emph{kennen}, \textquotedblleft{}to know\textquotedblright{} $\leftarrow$ Old High German \emph{kennan} $\leftarrow$ Proto-Germanic \emph{*kannijaną} (EtymDB 2.1): known.
\end{cousinweb}

\begin{grammarlens}[title={What you've built}]
One more word for this chapter.
\end{grammarlens}

\subsection*{Your turn}
\begin{itemize}
\raggedright
  \item \emph{Say it:} \emph{bekannt}
  \item \emph{Say it:} \emph{bekannt}, once more
  \item \emph{Say it:} say \emph{bekannt}
  \item \emph{Recall:} say the German for full, then the German for old-fashioned, then say \emph{bekannt} again
\end{itemize}

\begin{tcolorbox}[breakable,colback=teal!4,colframe=teal!35!black,title={Before you move on}]
What does bekannt mean? (\textquotedblleft{}Well known\textquotedblright{}.) Say it once more.
\end{tcolorbox}
\section[fremd]{fremd — foreign, strange}

\label{lesson:GE-C169-fremd}



\begin{quote}
Before the new one: say the German for old-fashioned, then the German for well known.
\end{quote}

\subsection*{You'll want to know: fremd}
\textbf{fremd} — \textquotedblleft{}foreign, strange\textquotedblright{}.

\begin{cousinweb}
\textbf{fremd} $\leftarrow$ Middle High German \emph{vremde}, from Old High German \emph{framadi}, \textquotedblleft{}away from here\textquotedblright{} — inherited; EtymDB 2.1 lists the old English \emph{fremd}, \textquotedblleft{}strange\textquotedblright{}, and Dutch \emph{vreemd} beside it.
\end{cousinweb}

\begin{grammarlens}[title={What you've built}]
One more word for this chapter.
\end{grammarlens}

\subsection*{Your turn}
\begin{itemize}
\raggedright
  \item \emph{Say it:} \emph{fremd}
  \item \emph{Say it:} \emph{fremd}, once more
  \item \emph{Say it:} say \emph{fremd}
  \item \emph{Recall:} say the German for old-fashioned, then the German for well known, then say \emph{fremd} again
\end{itemize}

\begin{tcolorbox}[breakable,colback=teal!4,colframe=teal!35!black,title={Before you move on}]
What does fremd mean? (\textquotedblleft{}Foreign, strange\textquotedblright{}.) Say it once more.
\end{tcolorbox}
\section[nervös]{nervös — nervous}

\label{lesson:GE-C169-nervoes}



\begin{quote}
Before the new one: say the German for well known, then the German for foreign.
\end{quote}

\subsection*{You'll want to know: nervös}
\textbf{nervös} — \textquotedblleft{}nervous\textquotedblright{}.

\begin{cousinweb}
\textbf{nervös} $\leftarrow$ Latin \emph{nervōsus}, \textquotedblleft{}sinewy, vigorous\textquotedblright{}, from \emph{nervus}, \textquotedblleft{}sinew, nerve\textquotedblright{} (EtymDB 2.1) — borrowed through French, like English \emph{nervous}.
\end{cousinweb}

\begin{grammarlens}[title={What you've built}]
One more word for this chapter.
\end{grammarlens}

\subsection*{Your turn}
\begin{itemize}
\raggedright
  \item \emph{Say it:} \emph{nervös}
  \item \emph{Say it:} \emph{nervös}, once more
  \item \emph{Say it:} say \emph{nervös}
  \item \emph{Recall:} say the German for well known, then the German for foreign, then say \emph{nervös} again
\end{itemize}

\begin{tcolorbox}[breakable,colback=teal!4,colframe=teal!35!black,title={Before you move on}]
What does nervös mean? (\textquotedblleft{}Nervous\textquotedblright{}.) Say it once more.
\end{tcolorbox}
\section[stolz]{stolz — proud}

\label{lesson:GE-C169-stolz}



\begin{quote}
Before the new one: say the German for foreign, then the German for nervous.
\end{quote}

\subsection*{You'll want to know: stolz}
\textbf{stolz} — \textquotedblleft{}proud\textquotedblright{}.

\begin{cousinweb}
\textbf{stolz} $\leftarrow$ Middle High German \emph{stolz} $\leftarrow$ Old High German \emph{stolz}, \textquotedblleft{}proud\textquotedblright{} (EtymDB 2.1). Where the Old High German word came from is uncertain.
\end{cousinweb}

\begin{grammarlens}[title={What you've built}]
One more word for this chapter.
\end{grammarlens}

\subsection*{Your turn}
\begin{itemize}
\raggedright
  \item \emph{Say it:} \emph{stolz}
  \item \emph{Say it:} \emph{stolz}, once more
  \item \emph{Say it:} say \emph{stolz}
  \item \emph{Recall:} say the German for foreign, then the German for nervous, then say \emph{stolz} again
\end{itemize}

\begin{tcolorbox}[breakable,colback=teal!4,colframe=teal!35!black,title={Before you move on}]
What does stolz mean? (\textquotedblleft{}Proud\textquotedblright{}.) Say it once more.
\end{tcolorbox}
\section[wütend]{wütend — angry}

\label{lesson:GE-C169-wuetend}



\begin{quote}
Before the new one: say the German for nervous, then the German for proud.
\end{quote}

\subsection*{You'll want to know: wütend}
\textbf{wütend} — \textquotedblleft{}angry\textquotedblright{}.

\begin{cousinweb}
\textbf{wütend} is the present participle of the verb \textquotedblleft{}to rage\textquotedblright{}, built on the word for \textquotedblleft{}fury\textquotedblright{}, Old High German \emph{wuot} $\leftarrow$ Proto-Germanic \emph{*wōdaz}, \textquotedblleft{}mad\textquotedblright{} (EtymDB 2.1). The god's name \emph{Wodan} shares the root.
\end{cousinweb}

\begin{grammarlens}[title={What you've built}]
One more word for this chapter.
\end{grammarlens}

\subsection*{Your turn}
\begin{itemize}
\raggedright
  \item \emph{Say it:} \emph{wütend}
  \item \emph{Say it:} \emph{wütend}, once more
  \item \emph{Say it:} say \emph{wütend}
  \item \emph{Recall:} say the German for nervous, then the German for proud, then say \emph{wütend} again
\end{itemize}

\begin{tcolorbox}[breakable,colback=teal!4,colframe=teal!35!black,title={Before you move on}]
What does wütend mean? (\textquotedblleft{}Angry\textquotedblright{}.) Say it once more.
\end{tcolorbox}
